\section{Introduction}

Inferring causal relations is a fundamental problem in the sciences and commercial applications. The problem of causal inference is often framed in terms of \emph{counterfactual} questions \citep{lewis1973causation,rubin1974estimating,pearl2009causality} such as ``Would this patient have lower blood sugar had she received a different medication?'', or ``Would the user have clicked on this ad had it been in a different color?''.
In this paper we propose a method to learn representations suited for counterfactual inference, and show its efficacy in both simulated and real world tasks.

We focus on counterfactual questions raised by what are known as \emph{observational studies}. Observational studies are studies where interventions and outcomes have been recorded, along with appropriate context. For example, consider an electronic health record dataset collected over several years, where for each patient we have lab tests and past diagnoses, as well as data relating to their diabetic status, and the causal question of interest is which of two existing anti-diabetic medications A or B is better for a given patient. Observational studies are rising in importance due to the widespread accumulation of data in fields such as healthcare, education, employment and ecology. We believe machine learning will be called on more and more to help make better decisions in these fields, and that researchers should be careful to pay attention to the ways in which these studies differ from classic supervised learning, as explained in Section \ref{Sec:prob} below.

In this work we draw a connection between counterfactual inference and domain adaptation. We then introduce a form of regularization by enforcing similarity between the distributions of representations learned for populations with different interventions. For example, the representations for patients who received medication A versus those who received medication B. This reduces the variance from fitting a model on one distribution and applying it to another. In Section \ref{Sec:model} we give several methods for learning such representations. In Section \ref{Sec:theory} we show our methods approximately minimizes an upper bound on a regret term in the counterfactual regime. The general method is outlined in Figure \ref{fig:model}.
Our work has commonalities with recent work on learning fair representations \citep{zemel2013learning,louizos2015variational} and learning representations for transfer learning \citep{ben2007analysis,gani2015domain}. In all these cases the learned representation has some invariance to specific aspects of the data: either an identity of a certain group such as racial minorities for fair representations, or the identity of the data source for domain adaptation, or, in the case of counterfactual learning, the type of intervention enacted in each population.

In machine learning, counterfactual questions typically arise in problems where there is a learning agent which performs actions, and receives feedback or reward for that choice without knowing what would be the feedback for other possible choices. This is sometimes referred to as bandit feedback \citep{beygelzimer2010contextual}. This setup comes up in diverse areas, for example off-policy evaluation in reinforcement learning \citep{sutton1998reinforcement}, learning from ``logged implicit exploration data'' \citep{strehl2010learning} or ``logged bandit feedback'' \citep{swaminathan2015batch}, and in understanding and designing complex real world ad-placement systems \cite{bottou2013counterfactual}.
Note that while in contextual bandit or robotics applications the researcher typically knows the method underlying the action choice (e.g. the policy in reinforcement learning), in observational studies we usually do not have control or even a full understanding of the mechanism which chooses which actions are performed and which feedback or reward is revealed. For instance, for anti-diabetic medication, more affluent patients might be insensitive to the price of a drug, while less affluent patients could bring this into account in their choice.

Given that we do not know beforehand the particulars determining the choice of action, the question remains, how can we learn from data which course of action would have better outcomes. By bringing together ideas from representation learning and domain adaptation, our method offers a novel way to leverage increasing computation power and the rise of large datasets to tackle consequential questions of causal inference.

The contributions of our paper are as follows. First, we show how to formulate the problem of counterfactual inference as a domain adaptation problem, and more specifically a covariate shift problem. Second, we derive new families of representation algorithms for counterfactual inference: one is based on linear models and variable selection, and the other is based on deep learning of representations \citep{bengio2013representation}. Finally, we show that learning representations that encourage similarity (balance) between the treated and control populations leads to better counterfactual inference; this is in contrast to many methods which attempt to create balance by re-weighting samples \citep[e.g.,][]{bang2005doubly,dudik2011doubly,austin2011introduction,swaminathan2015batch}. We show the merit of learning balanced representations both theoretically in Theorem 1, and empirically in a set of experiments across two datasets.


\section{Problem setup}
\label{Sec:prob}

Let $\mathcal{T}$ be the set of potential interventions or actions we wish to consider, $\mathcal{X}$ the set of contexts, and $\mathcal{Y}$ the set of possible outcomes. For example, for a patient $x \in \mathcal{X}$ the set $\mathcal{T}$ of interventions of interest might be two different treatments, and the set of outcomes might be $\mathcal{Y} = [0,200]$ indicating blood sugar levels in mg/dL. For an ad slot on a webpage $x$, the set of interventions $\mathcal{T}$ might be all possible ads in the inventory that fit that slot, while the potential outcomes could be $\mathcal{Y} = \{click, no\_click\}$. For a context $x$ (e.g. patient, webpage), and for each potential intervention $t \in \mathcal{T}$, let $Y_t(x)\in \mathcal{Y}$ be the \emph{potential outcome} for $x$. The fundamental problem of causal inference is that only one potential outcome is observed for a given context $x$: even if we give the patient one medication and later the other, the patient is not in exactly the same state. In machine learning this type of partial feedback is often called ``bandit feedback''. The model described above is known as the Rubin-Neyman causal model \citep{rubin1974estimating,rubin2011causal}.

We are interested in the case of a binary action set $\mathcal{T} = \{0,1\}$, where action $1$ is often known as the ``treated'' and action $0$ is the ``control''. In this case the quantity $Y_1(x) - Y_0(x)$ is of high interest: it is known as the \emph{individualized treatment effect} (ITE) for context $x$ \cite{van2007causal,weiss2015treatment}. Knowing this quantity enables choosing the best of the two actions when confronted with the choice, for example choosing the best treatment for a specific patient.
However, the fact that we only have access to the outcome of one of the two actions prevents the ITE from being known. Another commonly sought after quantity is the \emph{average treatment effect}, $\text{ATE} = \mathbb{E}_{x\sim p(x)}[\text{ITE}(x)]$ for a population with distribution $p(x)$. In the binary action setting, we refer to the observed and unobserved outcomes as the \emph{factual} outcome $y^F(x)$, and \emph{counterfactual} outcome $y^{CF}(x)$ respectively.

A common approach for estimating the ITE is by \emph{direct modelling}: given $n$ samples $\{(x_i,t_i,y^F_i)\}_{i=1}^n$, where $y^F_i = t_i \cdot Y_{1}(x_i) + (1-t_i) Y_{0}(x_i)$, learn a function $h: \mathcal{X} \times \mathcal{T} \rightarrow \mathcal{Y}$ such that $h(x_i,t_i) \approx y^F_i$. The estimated transductive ITE is then:
\begin{equation}
  \label{eq:ite}
\hat{\text{ITE}}(x_i) = \begin{cases}
        y^F_i - h(x_i,1-t_i), & t_i = 1.\\
        h(x_i,1-t_i) - y^F_i, & t_i = 0.
        \end{cases}
\end{equation}

While in principle any function fitting model might be used for estimating the ITE \citep{prentice1976use,gelman2006data,chipman2010bart,wager2015estimation,weiss2015treatment}, it is important to note how this task differs from standard supervised learning. The problem is as follows: the observed sample consists of the set $\hat{P}^F = \{(x_i,t_i)\}_{i=1}^n$. However, calculating the ITE requires inferring the outcome on the set $\hat{P}^{CF} = \{(x_i,1-t_i)\}_{i=1}^n$. We call the set $\hat{P}^F \sim P^F$ the empirical \emph{factual distribution}, and the set $\hat{P}^{CF} \sim P^{CF}$ the empirical \emph{counterfactual distribution}, respectively. Because $P^F$ and $P^{CF}$ need not be equal, the problem of causal inference by counterfactual prediction might require inference over a different distribution than the one from which samples are given. In machine learning terms, this means that the feature distribution of the test set differs from that of the train set. This is a case of \emph{covariate shift}, which is a special case of domain adaptation \citep{daume2006domain,jiang2008literature,mansour2009domain}. A somewhat similar connection was noted in \citet{ScholkopfJPSZMJ2012} with respect to covariate shift, in the context of a very simple causal model.

Specifically, we have that $P^F(x,t) = P(x) \cdot P(t|x)$ and $P^{CF}(x,t) = P(x) \cdot P(\neg t|x)$.  The difference between the observed (factual) sample and the sample we must perform inference on lies precisely in the treatment assignment mechanism, $P(t|x)$. For example, in a randomized control trial, we typically have that $t$ and $x$ are independent. In the contextual bandit setting, there is typically an algorithm which determines the choice of the action $t$ given the context $x$. In observational studies, which are the focus of this work, the treatment assignment mechanism is not under our control and in general will not be independent of the context $x$. Therefore, in general, the counterfactual distribution will be different from the factual distribution.

\begin{figure}[tbp!]
\centering
\includegraphics[width=0.95\columnwidth]{framework_color.pdf}
%\scalebox{1.0}{
%
\def\layersep{1.8cm}
\begin{tikzpicture}[->,>=latex,shorten >=1pt, thick,draw=black!50, node distance=\layersep]

    \tikzstyle{every pin edge}=[<-,shorten <=1pt]
    \tikzstyle{box}=[rectangle,minimum width=1.0cm,minimum height=1.0cm,draw=black!50];
    \tikzstyle{noline}=[box, fill=black!0, minimum width=0.6cm,minimum height=0.6cm, draw=black!0]];
    %\tikzstyle{input neuron}=[box, fill=green!15];

    % Draw the representation node
    \node[noline] (input) at (0,0) {$\bx$};
    \node[box, right of=input] (transform) {$g(\cdot)$};
    \node[noline, right of=transform] (rep) {$\phi$};
    \node[noline, below of=input] (group) {$t$};
    \node[noline, right of=rep] (output) {$\hat{\by}$};
    \node[noline, below of=output] (dist) {$\imb(\phi_C, \phi_T)$};
    %\node[hidden neuron] (lstm-1) at (\layersep,0) {$\bc_0$};
    %\node[output neuron, above of=lstm-1] (output-1) {$\bh_0$};

    \path (input) edge (transform);
    \path (transform) edge (rep);
    \path (rep) edge (output);
    \path (rep) edge (dist);
    \path (group) edge (output);
    \path (group) edge (dist);
    \path (group) edge (transform);
\end{tikzpicture}

%#}
\caption{\label{fig:model} Contexts $x$ are representated by $\Phi(x)$, which are used, with group indicator $t$, to predict the response $y$ while minimizing the imbalance in distributions measured by $\disc(\Phi_C, \Phi_T)$.}
\end{figure}

